\section{Prayers of the Season}\label{sec:prayers}

These are the prayers the book refers to, each printed from one named edition in its own spelling
and punctuation, except that the signs for versicle and response are printed V. and R. Where an English text is given, it is printed from the same edition as the Latin
or from the edition named under it; no translation has been made for this book.

\subsection{\latin{Rorate caeli}}

{\noindent\small PRAYER OF THE CHURCHES OF FRANCE DURING ADVENT. (Taken from the Prophet Isaias.)\par}

\begin{quote}
\textit{\latin{Rorate, coeli, desuper, et nubes pluant Justum.}}\par
\latin{Ne irascaris, Domine, ne ultra memineris iniquitatis: ecce civitas Sancti facta est deserta, Sion deserta facta est, Jerusalem desolata est, domus sanctificationis nostrae et gloriae tuae, ubi laudaverunt te patres nostri.}\par
\textit{\latin{Rorate, coeli, desuper, et nubes pluant Justum.}}\par
\latin{Peccavimus, et facti sumus tamquam immundus nos, et cecidimus quasi folium universi; et iniquitates nostrae quasi ventus abstulerunt nos: abscondisti faciem tuam a nobis, et allisisti nos in manu iniquitatis nostrae.}\par
\textit{\latin{Rorate, coeli, desuper, et nubes pluant Justum.}}\par
\latin{Vide, Domine, afflictionem populi tui, et emitte quem missurus es. Emitte Agnum dominatorem terrae de petra deserti ad montem filiae Sion, ut auferat ipse jugum captivitatis nostrae.}\par
\textit{\latin{Rorate, coeli, desuper, et nubes pluant Justum.}}\par
\latin{Consolamini, consolamini, popule meus: cito veniet salus tua: quare moerore consumeris? quare comprehendit te dolor? Salvabo te; noli timere: ego enim sum Dominus Deus tuus, Sanctus Israël, Redemptor tuus.}\par
\textit{\latin{Rorate, coeli, desuper, et nubes pluant Justum.}}\par
\end{quote}

\begin{quote}
\textit{Drop down dew, ye heavens, from above, and let the clouds rain the Just One.}\par
Be not angry, O Lord, and remember no longer our iniquity: behold the city of thy Sanctuary is become a desert, Sion is made a desert. Jerusalem is desolate, the house of our holiness and of thy glory, where our fathers praised thee.\par
\textit{Drop down dew, ye heavens, from above, and let the clouds rain the Just One.}\par
We have sinned, and we are become as one unclean, and we have all fallen as a leaf; and our iniquities, like the wind, have taken us away: thou hast hid thy face from us, and hast crushed us by the hand of our iniquity.\par
\textit{Drop down dew, ye heavens, from above, and let the clouds rain the Just One.}\par
See, O Lord, the affliction of thy people, and send Him whom thou hast promised to send. Send forth the Lamb, the ruler of the earth, from the rock of the desert to the mount of the daughter of Sion, that himself may take off the yoke of our captivity.\par
\textit{Drop down dew, ye heavens, from above, and let the clouds rain the Just One.}\par
Be comforted, be comforted, my people; thy salvation shall speedily come: why wilt thou waste away in sadness? why hath sorrow seized thee? I will save thee; fear not: for I am the Lord thy God, the Holy One of Israel, thy Redeemer.\par
\textit{Drop down dew, ye heavens, from above, and let the clouds rain the Just one.}\par
\end{quote}

{\smallskip\noindent\small Latin and English: Prosper Guéranger, \textit{The Liturgical Year: Advent}, tr. Laurence Shepherd (Dublin: Duffy, 1870), vol.~I, pp.~155--156. Guéranger gives it as a prayer of the Churches of France during Advent; its origin is not established. The last refrain reads ``the Just one'' as printed.\par}

\subsection{The Great O Antiphons}

At the Magnificat of Vespers, 17 to 23 December, each said entire before and after the canticle.

\noindent\textit{17 December.} \latin{O Sapientia, quae ex ore Altissimi prodiisti, attingens a fine usque ad finem, fortiter suaviterque disponens omnia: veni ad docendum nos viam prudentiae.}\par
\noindent O Wisdom, That comest out of the mouth of the Most High, That reachest from one end to another, and dost mightily and sweetly order all things: come, to teach us the way of prudence!\par
\smallskip

\noindent\textit{18 December.} \latin{O Adonai, et Dux domus Israël, qui Moysi in igne flammae rubi apparuisti, et ei in Sina legem dedisti: veni ad redimendum nos in brachio extento.}\par
\noindent O Adonai, and Ruler of the house of Israël, Who didst appear unto Moses in the burning bush, and gavest him the law in Sinai: come, to redeem us with an outstretched arm!\par
\smallskip

\noindent\textit{19 December.} \latin{O radix Jesse, qui stas in signum populorum, super quem continebunt reges os suum, quem Gentes deprecabuntur: veni ad liberandum nos, jam noli tardare.}\par
\noindent O Root of Jesse, Which standest for an ensign of the people, at Whom the kings shall shut their mouths, unto Whom the Gentiles shall seek: come, to deliver us, make no tarrying!\par
\smallskip

\noindent\textit{20 December.} \latin{O clavis David, et sceptrum domus Israël; qui aperis, et nemo claudit; claudis, et nemo aperit: veni, et educ vinctum de domo carceris, sedentem in tenebris, et umbra mortis.}\par
\noindent O Key of David, and Sceptre of the house of Israël; That openest, and no man shutteth; and shuttest and no man openeth: come, to bring out the prisoners from the prison, and them that sit in darkness, and in the shadow of death!\par
\smallskip

\noindent\textit{21 December.} \latin{O Oriens, splendor lucis aeternae, et sol justitiae: veni, et illumina sedentes in tenebris, et umbra mortis.}\par
\noindent O Day-Spring, Brightness of the everlasting Light, Sun of Righteousness, come, to give light to them that sit in darkness, and in the shadow of death!\par
\smallskip

\noindent\textit{22 December.} \latin{O Rex Gentium, et desideratus earum, lapisque angularis, qui facis utraque unum: veni, et salva hominem, quem de limo formasti.}\par
\noindent O King of the Gentiles, yea, and Desire thereof, O Corner-stone that makest of twain one: come, to save man, whom Thou hast made of the dust of the earth!\par
\smallskip

\noindent\textit{23 December.} \latin{O Emmanuel, Rex et legifer noster, exspectatio Gentium, et Salvator earum: veni ad salvandum nos Domine Deus noster.}\par
\noindent O Emmanuel, our King and our Law-giver, Longing of the Gentiles, yea, and Salvation thereof: come to save us, O Lord our God!\par
\smallskip


{\smallskip\noindent\small Latin: \textit{Breviarium Romanum}, pars hiemalis (Ratisbon: Pustet, 1888), pp.~156--157. English: tr. John, Marquess of Bute, \textit{The Roman Breviary}, Winter (Edinburgh and London: Blackwood, 1908), pp.~243--244.\par}

\subsection{\latin{Alma Redemptoris Mater}}

At the end of Compline, from Vespers of the Saturday before the First Sunday of Advent to the Purification; the versicle and prayer below are those of Advent.

\begin{verse}
\latin{Alma Redemptoris Mater, quae pervia coeli} \\
\latin{Porta manes, et stella maris, succurre cadenti,} \\
\latin{Surgere qui curat, populo: tu quae genuisti,} \\
\latin{Natura mirante, tuum sanctum Genitorem,} \\
\latin{Virgo prius ac posterius, Gabrielis ab ore} \\
\latin{Sumens illud Ave, peccatorum miserere.}
\end{verse}

\noindent\latin{V. Angelus Domini nuntiavit Mariae. R. Et concepit de Spiritu sancto.}\par
\noindent\latin{Oremus. Oratio.}\par
\noindent\latin{Gratiam tuam, quaesumus Domine, mentibus nostris infunde: ut qui, Angelo nuntiante, Christi Filii tui Incarnationem cognovimus, per passionem ejus et crucem ad resurrectionis gloriam perducamur. Per eumdem Christum Dominum nostrum. R. Amen.}\par
\medskip

\begin{verse}
Maiden! Mother of Him Who redeemed us, thou that abidest \\
Heaven's open gate, and the Star of the Sea, come, succour the fallen! \\
Fallen indeed we are, but fain would rise by thy succour. \\
Thou that beyond nature's course hast borne in time the Eternal; \\
Thou that a Virgin before, and after that childbirth remainest, \\
From the Archangel's lips the quickening message receiving, \\
Mother of JESUS and us, turn thine eyes of mercy on sinners.
\end{verse}

\noindent Verse. The Angel of the Lord announced unto Mary,\par
\noindent Answer. And she conceived of the Holy Ghost.\par
\noindent Let us pray.\par
\noindent We beseech Thee, O Lord, pour Thy grace into our hearts; that, as we have known the Incarnation of Thy Son JESUS Christ by the message of an Angel, so by His Passion and Cross we may be brought unto the glory of His Resurrection.\par
\noindent Through the same Christ our Lord.\par
\noindent Answer. Amen.\par

{\smallskip\noindent\small Latin: \textit{Breviarium Romanum}, pars hiemalis (Pustet, 1888), pp.~130--131. English: tr. Bute, \textit{The Roman Breviary}, Winter (1908), pp.~210--211, who in a footnote ascribes the antiphon to Hermann the Cripple (d.~1052). From First Vespers of Christmas both books give another versicle and prayer.\par}

\subsection{\latin{Creator alme siderum}}

The hymn at Vespers from the Saturday before the First Sunday of Advent.

\begin{verse}
\latin{Creator alme siderum,} \\
\latin{Aeterna lux credentium,} \\
\latin{Jesu Redemptor omnium,} \\
\latin{Intende votis supplicum.}

\latin{Qui daemonis ne fraudibus} \\
\latin{Periret orbis, impetu} \\
\latin{Amoris actus, languidi} \\
\latin{Mundi medela factus es.}

\latin{Commune qui mundi nefas} \\
\latin{Ut expiares; ad crucem} \\
\latin{E Virginis sacrario} \\
\latin{Intacta prodis victima.}

\latin{Cujus potestas gloriae,} \\
\latin{Nomenque cum primum sonat;} \\
\latin{Et coelites et inferi} \\
\latin{Tremente curvantur genu.}

\latin{Te deprecamur ultimae} \\
\latin{Magnum diei Judicem,} \\
\latin{Armis supernae gratiae} \\
\latin{Defende nos ab hostibus.}

\latin{Virtus, honor, laus, gloria} \\
\latin{Deo Patri cum Filio,} \\
\latin{Sancto simul Paraclito,} \\
\latin{In saeculorum saecula.} \\
\latin{Amen.}
\end{verse}

\begin{verse}
Creator of the starry pole, \\
Saviour of all who live, \\
And light of every faithful soul, \\
JESU, these prayers receive.

Who sooner than our foe malign \\
Should triumph, from above \\
Didst come, to be the medicine \\
Of a sick world, in love;

And the deep wounds to cleanse and cure \\
Of a whole race, didst go, \\
Pure Victim, from a Virgin pure, \\
The bitter Cross unto.

Who hast a Name, and hast a Power, \\
The height and depth to sway, \\
And Angels bow, and devils cower, \\
In transport or dismay;

Thou too shalt be our Judge at length; \\
Lord, in Thy grace bestow \\
Thy weapons of celestial strength, \\
And snatch us from the foe.

Honour and glory, power and praise, \\
To Father, and to Son, \\
And Holy Ghost, be paid always, \\
The Eternal Three in One. Amen.
\end{verse}

\noindent\latin{V. Rorate coeli desuper, et nubes pluant justum. R. Aperiatur terra, et germinet Salvatorem.}\par
\noindent Verse. Drop down, ye heavens, from above, and let the skies pour down the Righteous One.\par
\noindent Answer. Let the earth open, and let her bring forth the Saviour.\par

{\smallskip\noindent\small Latin hymn and versicle: \textit{Breviarium Romanum}, pars hiemalis (Pustet, 1888), p.~132. English: the translation of Cardinal Newman, as Bute prints and attributes it, \textit{The Roman Breviary}, Winter (1908), p.~213, with Bute's English of the versicle.\par}

\subsection{The Angelus}

\noindent\latin{V. Angelus DOMINI nuntiavit Mariae.}\par
\noindent\latin{R. Et concepit de SPIRITU SANCTO. Ave Maria.}\par
\noindent\latin{V. Ecce ancilla DOMINI.}\par
\noindent\latin{R. Fiat mihi secundum verbum tuum. Ave Maria.}\par
\noindent\latin{V. Et Verbum caro factum est.}\par
\noindent\latin{R. Et habitavit in nobis. Ave Maria.}\par
\noindent\latin{V. Ora pro nobis, sancta DEI Genitrix.}\par
\noindent\latin{R. Ut digni efficiamur promissionibus CHRISTI.}\par
\noindent\latin{Oremus.}\par
\noindent\latin{Gratiam tuam, quaesumus, DOMINE, mentibus nostris infunde: ut qui, angelo nuntiante, CHRISTI FILII tui Incarnationem cognovimus, per Passionem ejus et Crucem ad resurrectionis gloriam perducamur. Per eundem CHRISTUM DOMINUM nostrum. R. Amen.}\par
\medskip

\noindent V. The angel of the LORD declared unto Mary.\par
\noindent R. And she conceived of the HOLY GHOST. Ave Maria.\par
\noindent V. Behold the handmaid of the LORD.\par
\noindent R. Be it done unto me according to thy word. Ave Maria.\par
\noindent V. And the Word was made Flesh.\par
\noindent R. And dwelt among us. Ave Maria.\par
\noindent V. Pray for us, O Holy Mother of GOD.\par
\noindent R. That we may be made worthy of the promises of CHRIST.\par
\noindent Let us pray.\par
\noindent Pour forth, we beseech Thee, O LORD, thy grace into our hearts; that we, to whom the Incarnation of CHRIST thy SON was made known by the message of an angel, may by his Passion and Cross be brought to the glory of his resurrection. Through the same CHRIST our LORD. R. Amen.\par

{\smallskip\noindent\small Latin and English: \textit{The Raccolta}, tr. Ambrose St John (London: Burns and Oates, 1910), no.~182, pp.~151--152, with its capitals. The Raccolta does not print the Hail Mary at each \latin{Ave Maria}.\par}

\subsection{The St Andrew Prayer}

\begin{quote}
ASPIRATION.\par
Hail and blessed be the hour and moment in which the Son of God was born of the most pure Virgin Mary at midnight in Bethlehem in the piercing cold. In that hour vouchsafe, O my God! to hear my prayer and grant my desires, through the merits of Our Saviour Jesus Christ, and of His blessed Mother. Amen.\par
It is piously believed that whoever recites the above prayer fifteen times a day, from the feast of St. Andrew (Nov. 30) until Christmas, will obtain what is asked.\par
\end{quote}

{\smallskip\noindent\small A. A. Lings, \textit{Our Favorite Novenas} (New York: Benziger, 1897), p.~29, where it is printed as an ``Aspiration'' without a name. The note printed under it records a pious belief and makes no promise. The book carries the \latin{Nihil obstat} of Thomas L. Kinkead and the \latin{Imprimatur} of Michael Augustine, Archbishop of New York, 6 February 1897, which approve the compilation.\par}

\subsection{An Act of Contrition}

\begin{quote}
O my God! I am heartily sorry for having offended Thee, and I detest all my sins, because I dread the loss of heaven and the pains of hell, but most of all because they offend Thee, my God, who art all-good and deserving of all my love. I firmly resolve, with the help of Thy grace, to confess my sins, to do penance, and to amend my life. Amen.\par
\end{quote}

{\smallskip\noindent\small \textit{A Manual of Prayers for the Use of the Catholic Laity} (New York: Catholic Publication Society, 1889), p.~54. The same book at p.~286, and the Baltimore Catechism, print slightly different forms; this one is printed alone.\par}

